% GENERATED FILE. Edit canonical lessons, then run npm run generate:books.
% canonical-source-hash: fnv1a64:249def80e3b2889b
% canonical-lessons: GU-C317-itihas, GU-C317-sadi, GU-C317-pracin, GU-C317-bhutkal, GU-C317-janmavu

\chapter{Pages of History}
\label{ch:gu-pages-of-history}
\hlchaptermodality{317}

\begin{chapteropening}
\textbf{By the end of this chapter:} \emph{I can say when and where something old began.}

Complete the last lesson of chapter 317: I can say when and where something old began.
\end{chapteropening}

\section[\texorpdfstring{\gu{ઇતિહાસ} (itihās)}{ઇતિહાસ (itihās)}]{\gu{ઇતિહાસ} (itihās) — history}

\label{lesson:GU-C317-itihas}



\begin{quote}
Before the new one: say the Gujarati for secret, then the Gujarati for to last.
\end{quote}

\subsection*{You'll want to know: \gu{ઇતિહાસ}}
\textbf{\gu{ઇતિહાસ}} — \emph{itihās} — \textquotedblleft{}history\textquotedblright{}.

\begin{quote}
\gu{મીરાને} \gu{ઇતિહાસ} \gu{ખૂબ} \gu{ગમે} \gu{છે}.
\end{quote}

\begin{quote}
\emph{Mīrāne itihās khūb game che.}
\end{quote}

\begin{quote}
Mira loves history.
\end{quote}

A masculine noun.

\begin{cousinweb}
\textbf{\gu{ઇતિહાસ}} (\emph{itihās}) is a tatsama, Sanskrit \emph{itihāsa}, from \emph{iti ha āsa}, \textquotedblleft{}so indeed it was\textquotedblright{} (Monier-Williams): the telling of what happened long ago. It is masculine: \textbf{\gu{ગુજરાતનો} \gu{ઇતિહાસ}}, the history of Gujarat.
\end{cousinweb}

\begin{grammarlens}[title={What you've built}]
One more word for telling what happened.
\end{grammarlens}

\subsection*{Your turn}
\begin{itemize}
\raggedright
  \item \emph{Say it:} \emph{itihās}
  \item \emph{Say it:} \emph{itihās}, once more
  \item \emph{Say it:} \emph{Mīrāne itihās khūb game che.} — Mira loves history.
  \item \emph{Recall:} say the Gujarati for secret, then the Gujarati for to last, then say \emph{itihās} again
\end{itemize}

\begin{tcolorbox}[breakable,colback=teal!4,colframe=teal!35!black,title={Before you move on}]
What does \gu{ઇતિહાસ} mean? (\textquotedblleft{}History\textquotedblright{}.) Say it once more.
\end{tcolorbox}
\section[\texorpdfstring{\gu{સદી} (sadī)}{સદી (sadī)}]{\gu{સદી} (sadī) — a century}

\label{lesson:GU-C317-sadi}



\begin{quote}
Before the new one: say the Gujarati for to last, then the Gujarati for history.
\end{quote}

\subsection*{You'll want to know: \gu{સદી}}
\textbf{\gu{સદી}} — \emph{sadī} — \textquotedblleft{}a century\textquotedblright{}.

\begin{quote}
\gu{આ} \gu{મંદિર} \gu{ત્રણ} \gu{સદી} \gu{જૂનું} \gu{છે}.
\end{quote}

\begin{quote}
\emph{ā mandir traṇ sadī jūnũ che.}
\end{quote}

\begin{quote}
This temple is three centuries old.
\end{quote}

A feminine noun.

\begin{cousinweb}
\textbf{\gu{સદી}} (\emph{sadī}) is usually traced to Persian \emph{sadī}, \textquotedblleft{}a century\textquotedblright{}, built on the word for a hundred. It is feminine: \textbf{\gu{આ} \gu{સદીમાં}}, in this century.
\end{cousinweb}

\begin{grammarlens}[title={What you've built}]
One more word for telling what happened.
\end{grammarlens}

\subsection*{Your turn}
\begin{itemize}
\raggedright
  \item \emph{Say it:} \emph{sadī}
  \item \emph{Say it:} \emph{sadī}, once more
  \item \emph{Say it:} \emph{ā mandir traṇ sadī jūnũ che.} — This temple is three centuries old.
  \item \emph{Recall:} say the Gujarati for to last, then the Gujarati for history, then say \emph{sadī} again
\end{itemize}

\begin{tcolorbox}[breakable,colback=teal!4,colframe=teal!35!black,title={Before you move on}]
What does \gu{સદી} mean? (\textquotedblleft{}A century\textquotedblright{}.) Say it once more.
\end{tcolorbox}
\section[\texorpdfstring{\gu{પ્રાચીન} (prācīn)}{પ્રાચીન (prācīn)}]{\gu{પ્રાચીન} (prācīn) — ancient}

\label{lesson:GU-C317-pracin}



\begin{quote}
Before the new one: say the Gujarati for history, then the Gujarati for a century.
\end{quote}

\subsection*{You'll want to know: \gu{પ્રાચીન}}
\textbf{\gu{પ્રાચીન}} — \emph{prācīn} — \textquotedblleft{}ancient\textquotedblright{}.

\begin{quote}
\gu{આ} \gu{ગામમાં} \gu{એક} \gu{પ્રાચીન} \gu{કૂવો} \gu{છે}.
\end{quote}

\begin{quote}
\emph{ā gāmmā̃ ek prācīn kūvo che.}
\end{quote}

\begin{quote}
There is an ancient well in this village.
\end{quote}

\begin{cousinweb}
\textbf{\gu{પ્રાચીન}} (\emph{prācīn}) is a tatsama, Sanskrit \emph{prācīna}, \textquotedblleft{}turned to the east\textquotedblright{}, and so \textquotedblleft{}former, ancient\textquotedblright{} (Monier-Williams). It is for things of long ago; \textbf{\gu{જૂનું}} is old in the everyday sense.
\end{cousinweb}

\begin{grammarlens}[title={What you've built}]
One more word for telling what happened.
\end{grammarlens}

\subsection*{Your turn}
\begin{itemize}
\raggedright
  \item \emph{Say it:} \emph{prācīn}
  \item \emph{Say it:} \emph{prācīn}, once more
  \item \emph{Say it:} \emph{ā gāmmā̃ ek prācīn kūvo che.} — There is an ancient well in this village.
  \item \emph{Recall:} say the Gujarati for history, then the Gujarati for a century, then say \emph{prācīn} again
\end{itemize}

\begin{tcolorbox}[breakable,colback=teal!4,colframe=teal!35!black,title={Before you move on}]
What does \gu{પ્રાચીન} mean? (\textquotedblleft{}Ancient\textquotedblright{}.) Say it once more.
\end{tcolorbox}
\section[\texorpdfstring{\gu{ભૂતકાળ} (bhūtkāḷ)}{ભૂતકાળ (bhūtkāḷ)}]{\gu{ભૂતકાળ} (bhūtkāḷ) — the past}

\label{lesson:GU-C317-bhutkal}



\begin{quote}
Before the new one: say the Gujarati for a century, then the Gujarati for ancient.
\end{quote}

\subsection*{You'll want to know: \gu{ભૂતકાળ}}
\textbf{\gu{ભૂતકાળ}} — \emph{bhūtkāḷ} — \textquotedblleft{}the past\textquotedblright{}.

\begin{quote}
\gu{ભૂતકાળમાં} \gu{આ} \gu{શહેર} \gu{ખૂબ} \gu{નાનું} \gu{હતું}.
\end{quote}

\begin{quote}
\emph{bhūtkāḷmā̃ ā śaher khūb nānũ hatũ.}
\end{quote}

\begin{quote}
In the past this city was very small.
\end{quote}

A masculine noun.

\begin{cousinweb}
\textbf{\gu{ભૂતકાળ}} (\emph{bhūtkāḷ}) is built from Sanskrit \emph{bhūta}, \textquotedblleft{}been, gone, past\textquotedblright{}, and \emph{kāla}, \textquotedblleft{}time\textquotedblright{} (Monier-Williams): the time that has been. It is masculine: \textbf{\gu{ભૂતકાળમાં}}, in the past.
\end{cousinweb}

\begin{grammarlens}[title={What you've built}]
One more word for telling what happened.
\end{grammarlens}

\subsection*{Your turn}
\begin{itemize}
\raggedright
  \item \emph{Say it:} \emph{bhūtkāḷ}
  \item \emph{Say it:} \emph{bhūtkāḷ}, once more
  \item \emph{Say it:} \emph{bhūtkāḷmā̃ ā śaher khūb nānũ hatũ.} — In the past this city was very small.
  \item \emph{Recall:} say the Gujarati for a century, then the Gujarati for ancient, then say \emph{bhūtkāḷ} again
\end{itemize}

\begin{tcolorbox}[breakable,colback=teal!4,colframe=teal!35!black,title={Before you move on}]
What does \gu{ભૂતકાળ} mean? (\textquotedblleft{}The past\textquotedblright{}.) Say it once more.
\end{tcolorbox}
\section[\texorpdfstring{\gu{જન્મવું} (janmavũ)}{જન્મવું (janmavũ)}]{\gu{જન્મવું} (janmavũ) — to be born}

\label{lesson:GU-C317-janmavu}



\begin{quote}
Before the new one: say the Gujarati for ancient, then the Gujarati for the past.
\end{quote}

\subsection*{You'll want to know: \gu{જન્મવું}}
\textbf{\gu{જન્મવું}} — \emph{janmavũ} — \textquotedblleft{}to be born\textquotedblright{}.

\begin{quote}
\gu{મારા} \gu{દાદા} \gu{આ} \gu{ગામમાં} \gu{જન્મ્યા} \gu{હતા}.
\end{quote}

\begin{quote}
\emph{mārā dādā ā gāmmā̃ janmyā hatā.}
\end{quote}

\begin{quote}
My grandfather was born in this village.
\end{quote}

Past (he …): \gu{તે} \gu{અહીં} \gu{જન્મ્યો}, \emph{te ahī̃ janmyo}.

\begin{cousinweb}
\textbf{\gu{જન્મવું}} (\emph{janmavũ}) is built on \emph{janma}, birth, from Sanskrit \emph{jánman}, \textquotedblleft{}birth\textquotedblright{} (CDIAL 5113), the word in \textbf{\gu{જન્મદિવસ}}, a birthday. The past is plural for an elder, out of respect: \textbf{\gu{દાદા} \gu{જન્મ્યા} \gu{હતા}}.
\end{cousinweb}

\begin{grammarlens}[title={What you've built}]
One more word for telling what happened.
\end{grammarlens}

\subsection*{Your turn}
\begin{itemize}
\raggedright
  \item \emph{Say it:} \emph{janmavũ}
  \item \emph{Say it:} \emph{janmavũ}, once more
  \item \emph{Say it:} \emph{mārā dādā ā gāmmā̃ janmyā hatā.} — My grandfather was born in this village.
  \item \emph{Recall:} say the Gujarati for ancient, then the Gujarati for the past, then say \emph{janmavũ} again
\end{itemize}

\begin{tcolorbox}[breakable,colback=teal!4,colframe=teal!35!black,title={Before you move on}]
What does \gu{જન્મવું} mean? (\textquotedblleft{}To be born\textquotedblright{}.) Say it once more.
\end{tcolorbox}
